% Figure: mean squared error of two estimators of a uniform endpoint,
% the moment estimator 2*Xbar (unbiased, MSE = theta^2/(3n)) and the
% sample maximum (biased, MSE = 2 theta^2 /((n+1)(n+2))), plotted in
% units of theta^2 against sample size n. The maximum wins for every
% n and its advantage widens as 1/n^2 against the moment estimator's
% 1/n. A third curve, the bias-corrected maximum with MSE
% theta^2/(n(n+2)), sits lowest of all.
\begin{figure}[htbp]
\centering
\begin{tikzpicture}
\begin{axis}[
    width=12cm, height=7cm,
    xmin=2, xmax=20,
    ymin=0, ymax=0.18,
    axis lines=left,
    xlabel={Sample size $n$},
    ylabel={Mean squared error (units of $\theta^2$)},
    every axis plot/.append style={very thick, mark=none},
    legend style={font=\small, draw=none, fill=none, at={(0.98,0.95)}, anchor=north east},
    xtick={2,4,6,8,10,12,14,16,18,20},
    domain=2:20, samples=60,
]
% Moment estimator 2*Xbar: MSE = 1/(3n).
\addplot[engblue]{1/(3*x)};
\addlegendentry{moment estimator $2\bar{X}$}
% Sample maximum: MSE = 2/((n+1)(n+2)).
\addplot[engred]{2/((x+1)*(x+2))};
\addlegendentry{sample maximum $\max_i X_i$}
% Bias-corrected maximum: MSE = 1/(n(n+2)).
\addplot[englime!70!black]{1/(x*(x+2))};
\addlegendentry{bias-corrected maximum}
\end{axis}
\end{tikzpicture}
\caption[Mean squared error of two uniform-endpoint estimators]{Mean
squared error, in units of $\theta^2$, for three estimators of the
upper endpoint of a $\mathrm{Uniform}[0,\theta]$ population. The
unbiased moment estimator $2\bar{X}$ falls only as $1/n$, while the
biased sample maximum falls as $1/n^2$ and beats the moment estimator
for every sample size. The bias-corrected maximum $\frac{n+1}{n}\max_i
X_i$ removes the bias and sits lowest of all. The picture is the
chapter's standing argument that unbiasedness is not the objective and
mean squared error is: a biased estimator with a much smaller variance
is the better report.}
\label{fig:vol02:ch14:mse-tradeoff-uniform}
\end{figure}
